% ===========================================================================
%  PARKED --- not compiled, not deleted.
%
%  Source: CH9/document MAIN/sections/07_bardo_thodol.tex, lines 1-53 (the
%  whole section), together with its three figure inclusions
%  (IMG_ch8/samsara.jpg, IMG_ch8/proverbe.jpg, root/lightning.jpg).  The image
%  files themselves are in ../figures/, flattened and with spaces in file
%  names replaced by underscores, so this file can be restored by changing the
%  three \includegraphics paths to samsara, proverbe and lightning.
%
%  Why parked: confirmed cut in CH9_DECISIONS.md section 4.  It belongs to no
%  model.  Verified before parking, per plan section 4.1, that nothing
%  mathematical is lost:
%    - the one substantive idea, that near a critical point small
%      perturbations are amplified, is made quantitatively in this chapter's
%      Pimentel section at \cref{var:eq:icrit} and in the prose around it;
%    - quasi-stationarity is treated properly in the distribution-theory
%      chapter (dist:sec:qsd, dist:thm:qsd, dist:eq:qsd, dist:sec:qsd-geometric)
%      and defined in the mathematical introduction (m:def:qsd, m:eq:qsmean).
%  Both were checked against those chapters' label lists on 21 August 2026.
%
%  The three multiply-defined \label{redone} instances below are part of the
%  defect this rewrite repairs; they are left here exactly as found.
% ===========================================================================

\section*{Bardo Thodol}

 In the Bardo Thodol (otherwise known as The Tibetan book of living and dying), there are listed 6 bardo states which describe the stages of life and the process of transition between lives:

\begin{itemize} \item the bardo of birth \item the bardo of dreams \item the bardo of samadhi meditation \item the bardo of the moment before death \item the bardo of dharmata (pure reality) \item the bardo of becoming \end{itemize}

I can't help but notice a crossover with the first three of these and the stages of the Mandukya Upanishad: A, U, and M for "beginning," "process," and "cessation." This links loosely with the stages of the logistic growth model. Initially, the number of species rapidly increases, then there is a stabilizing phase, and eventually, many die out, leaving only a few. But this is an aside.

The point of the Bardo Thodol demonstrates an understanding of something also revealed by stochastic process theory: that in life, there are periods of relative stasis punctuated by times of great upheaval, change, and above all, uncertainty. The book is a guide for those dying, so that they may pass through the crucial final three stages to attain a favorable passage through the period of undetermined futures. The idea is that their consciousness somehow persists after their physical body dies. It then wanders through the bardo of dharmata, before stumbling upon a series of lights. Each light encountered corresponds to one of the 6 realms of samsara in which the person might find themselves born.

\vspace{5pt} \hrule \vspace{5pt}

\noindent \textbf{Aside -- the 6 realms of samsara: } In Tibetan Buddhist theory, there are 6 "realms" in which all beings supposedly occupy, cycling among them over a series of cycles of birth, death, and rebirth:

\begin{itemize} \item Hell realm(s) \item Animal realm \item Realm of hungry ghosts \item Human realm \item Realm of the jealous gods \item Deva realm \end{itemize}

The stated ideal is to, as it were, exit this process of cycles and become "enlightened." This can only happen at the time of death and supposedly is most easily accomplished when exiting the human realm. The trick is to use the opportunity of non-fixation to remain in the state of transience without falling back into one of the realms, and it is this which the Bardo Thodol aims to guide its recipients towards.

\vspace{5pt} \hrule \vspace{5pt}

The lights seem to be analogous to circumstances that seem appealing but end up as "basins of attraction," whereby entry inevitably leads to an absorption and fixation in that particular realm. There is a clear analogy to observations in stochastic process analysis. For example, at a state of undetermined instability (as in the squeezes that take place when one population is very small in the Pimental model), nearby trajectories diverge as a result of an amplification in the importance of small, otherwise insignificant, stochastic influences.

In this framework, when a person dies, their fated trajectory is yet undecided, and in this state of undetermined instability, small influences can be applied to steer things one way or the other.

\subsubsection*{Why do we care?} In the modern scientific age we live in, most educated people balk and sneer at the concept of life after death.  There appears two angles through which to address the matter:

\begin{itemize} \item \textbf{You could say that there are many lives within one life}, and that the transition between phases of life shares qualities with the stated transition between the bardo of the moment before death and the bardo of becoming. In a way, I feel as if I am embarking on the bardo of the moment before death.  The path of life after a PhD, without concrete plans, can be underermined and subject to an array of possible continuations. 

\item (More abstractly) \textbf{Multiple successive individual people can occupy the same character.} It can be imagined that this effect will have been all the more pronounced in 9th-century Tibet. For centuries, over countless generations, families will have occupied the Tibetan plateau in relative isolation. With this slowed movement, it seems only natural that the perception of rise, fall, and change will have been extended to perceive the agencies living over multiple generations. Take as a prime example the "character" of the Dalai Lama and the continuously reincarnated lamas of the other major schools: one personality, living over 10 or more generations.

New head lamas are selected at birth. They are identified by the elders of the school to be the "reincarnation" of the previous leader. It strikes me as a remarkable mechanism of government for creating piety and meritous behaviour amongst a people: give everyone's child the prospect of becoming leader of the state.

\end{itemize}

\subsection*{Quasistationarity} 

Verging a bit into the reckless and speculative here.  But is there not a philosophical analogy between quasistationarity --- that is,  quantified behaviour of a stochastic process given non-fixation --- and the theoretical description of liberation in the conception of the Bardo Thodol. \par 
Not getting stuck in one of the basins of attraction of the 6 realms, appears similar to not getting stuck, either by depletion to 0 or by catastrophic fixation in the birth death catastrophe process.  Similar to Icarus, neither flying too high or too low,  the liberatee of the bardo of dharmata doesn't get too close to any of the appealing absorbative attractions which would inevitably lead to a re-binding into one of the cyclic domains of samsara.  \par 
Behavior conditional on non-fixation, non rebirth --- that being the practive of the Arahant as that of the Bodisatva.  Unbound and unmarred by conditions of suffering which arise from the succeptability of craving for material manifestations. \par 



\begin{figure}[H] \includegraphics[width = 0.85\textwidth]{figures/IMG_ch8/samsara.jpg}
\centering
\pagebreak \caption{Wheel of samsara. Yama (god of death) depicted holding a wheel containing the 6 realms. From the center, a rooster, snake, and pig are shown to represent attachment, aversion, and ignorance respectively. They are said to link with one another cyclically, causing the dynamism which makes the wheel progress.} \label{redone} \end{figure}

\begin{figure}[H] \includegraphics[width = \textwidth]{figures/IMG_ch8/proverbe} \caption{Pieter Bruegel’s Netherlandish Proverbs (oil on wood, painted 1559). For years, this painting has fascinated me. One of the reasons is that it was painted in a time still marred by plagues. Enough time had persisted since the first wave in 1345 for the societal upheaval to become fully manifest. Above all, it appears to depict a time in which the previous structures, customs, and societal norms are in complete disarray. A time of flux and dynamism, perhaps representative of what Nietzsche said in the quote referenced above: \ \noindent"At this turning-point of history there manifest themselves, side by side, and often mixed and entangled together, a magnificent, manifold, virgin-forest-like up-growth and up-striving, a kind of tropical tempo in the rivalry of growth, and an extraordinary decay and self-destruction, owing to the savagely opposing and seemingly exploding egoisms, which strive with one another 'for sun and light.'"} \label{redone} \end{figure}



\begin{figure}[H] \includegraphics[width = \textwidth]{figures/root/lightning}
\centering
\pagebreak \caption{Coach in a Thunderstorm (Philippe Jacques de Loutherbourg, 1790).  It almost appears as if the cart, the horse and the person are eminating from the lightning strike,  rather than merely running away from it after being struck on the road. As if the energy of the lightning itself is creating the activity we see before us.  We have no way of knowing whether Philippe Jacques intended for this this interpretation.  But the emergences of sudden complex order through immediate resolutions of large potential,  seems like it might be consistent with the situaten of  a country (France) at the verge of instability,  which ultimately manifested into full blown societal revolution only a few years later.  To be clear, I am suggesting the interpretation that the motion depected is an ephemeral dissipative structure,  momentarily created by the lightning. } \label{redone} \end{figure}
